\section{Second-row stabilization and actual multiplicities}

\subsection*{Finding F-014: stabilization controlled by the second row}
Let \(d\geq1\), \(n\geq1\), and let \(\lambda\vdash nd\)
have \(r\leq d\) rows. If \(r=1\), its ambient multiplicity is
one and it occurs in the determinant coordinate ring. Suppose
\(r\geq2\), put \(b=\lambda_2\), and assume \(n\geq b\).
Then
\[
       \mu=(db-|\bar\lambda|,\lambda_2,\ldots,\lambda_r)
                                                     \vdash db
\]
is a partition and
\begin{equation}\label{eq:F014stab}
                         a_\lambda(d[n])=a_\mu(d[b]).
\end{equation}
More precisely, multiplication by the fixed first basis vector
in each symmetric-power factor identifies the entire highest-weight
spaces of these two weights. The statement is established here
directly from Cauchy and flag minors; no assertion of historical
priority for plethysm stabilization is intended.

\begin{proof}
Let \(U\) have dimension \(r\), let \(E=\mathbb C^d\), and use
variables \(z_{ij}\), \(1\leq i\leq r\), \(1\leq j\leq d\),
for the symmetric algebra \(\operatorname{Sym}(U\otimes E)\).
Rows carry the \(\mathrm{GL}(U)\) action and columns the commuting
\(\mathrm{GL}(E)\) action. The Cauchy decomposition says that
the space \(H_\lambda\) of row highest-weight vectors of weight
\(\lambda\) is the irreducible column module \(S_\lambda E\).
The condition \(r\leq d\) ensures that this module is nonzero.

For a column subset \(J\) of size \(h\), call
\[
             \Delta_{h,J}=\det(z_{ij})_{1\leq i\leq h,\ j\in J}
\]
a flag minor. It is a row highest-weight vector of weight
\((1^h)\). Consider all products of flag minors whose heights,
with multiplicity, are the column lengths \(\lambda'_j\).
Each product is a row highest-weight vector of weight \(\lambda\).
Their linear span is stable under \(\mathrm{GL}(E)\), because
the column change-of-basis formula for a minor expresses it
as a linear combination of minors of the same height.
This span is nonzero: choose the first \(h\) columns for every
minor of height \(h\); their product is a nonzero polynomial.
Irreducibility of \(H_\lambda\simeq S_\lambda E\) now proves
that these products span all of \(H_\lambda\).
This proves the required spanning assertion without presuming
an unproved basis theorem for the products.

Exactly \(\lambda_2=b\) of the minor factors have height at
least two. In the expansion of any one minor, each column is
used at most once. Height-one minors use only first-row variables.
It follows that every monomial in every spanning product, and
therefore in every element of \(H_\lambda\), has degree at most
\(b\) in the variables below the first row of any fixed column.

Take the column multidegree \((n,\ldots,n)\) part of
\(H_\lambda\). Every monomial there consequently has exponent
at least \(n-b\) on each \(z_{1j}\). Hence every polynomial in
that space is divisible by
\[
                         Q^{n-b},\qquad Q=\prod_{j=1}^d z_{1j}.
\]
The polynomial \(Q\) is fixed by the upper row unipotent group,
has row weight \(de_1\), and is fixed by column permutations.
Division by \(Q^{n-b}\) therefore produces a row highest-weight
vector of weight \(\mu=\lambda-d(n-b)e_1\), of column
multidegree \((b,\ldots,b)\). Conversely multiplication by
\(Q^{n-b}\) sends every such vector back to the original space.
These operations are inverse and commute with column permutations.

The tuple \(\mu\) is a partition because
\[
 |\bar\lambda|\leq(r-1)b\leq(d-1)b,
                  \qquad db-|\bar\lambda|\geq b=\lambda_2.
\]
Finally, column-permutation invariants in the balanced column
space are \(\operatorname{Sym}^d\operatorname{Sym}^n U\)
and \(\operatorname{Sym}^d\operatorname{Sym}^b U\), respectively.
Taking invariants in the displayed isomorphism proves
\eqref{eq:F014stab}, including all multiplicity copies.
The one-row case consists of a single highest-weight monomial
and follows directly.
\end{proof}

\paragraph{Full determinant multiplicities.}
For \(b\geq1\), \(r\geq2\), write
\[
 u=\left\lfloor\frac{b-1}{2}\right\rfloor,\qquad
 v=\left\lceil\frac{b-1}{2}\right\rceil,\qquad
 B(b,r)=\binom{r-1+u}{u}+\binom{r-1+v}{v}-1.
\]
If \(r\leq n^2\) and \(n\geq B(b,r)\), then
\begin{equation}\label{eq:F014det}
                         K_n(\lambda)=a_\lambda(d[n]).
\end{equation}
In particular there is no actual permanent-versus-determinant
multiplicity obstruction of this shape, at any outer degree.

\begin{proof}
The degree-splitting construction O-012.a puts every form
\(x_1^{n-b}f\), \(f\in\operatorname{Sym}^b U^*\), in
\(\Omega_n\) under the stated inequality.
It also gives \(B(b,r)\geq b\), so the stabilization applies.
By the arbitrary-rank padding inequality O-012.b,
\[
 K_n(\mu+d(n-b)e_1)\geq a_\mu(d[b])=a_\lambda(d[n]).
\]
Since the partition on the left is \(\lambda\), and an orbit-closure
coordinate ring is an ambient quotient, the reverse inequality
also holds. This proves \eqref{eq:F014det}.

The distinction between the two maps used here is essential.
Multiplication by \(e_1^{n-b}\) on the coordinate tensor space
is Borel equivariant up to a character and gives the highest-weight
space isomorphism. Actual pullback to padded forms is its
differential transpose. Their composition is diagonal with
nonzero factorial entries, as proved in O-012.b; it preserves
linear independence of every lifted copy. No false assertion
that those two maps are the same has been used.
\end{proof}

\paragraph{Explicit GCT second-row exclusions.}
For \(m\geq2\), a positive permanent-side multiplicity satisfies
\(r\leq m^2\) in the internal-padding convention.
The number \(\binom{r-1+j}{j}\) counts degree-\(j\) monomials
in \(r\) variables, so it is at most \(r^j\), by mapping
words of length \(j\) to their commutative monomials.
Thus
\[
                         B(b,r)\leq r^u+r^v-1.
\]
If \(b\leq5\), then \(u,v\leq2\), and
\[
               B(b,r)\leq2r^2-1\leq2m^4-1<3m^4.
\]
Consequently in the range \(n>3m^4\), every actual multiplicity
obstruction must satisfy
\begin{equation}\label{eq:F014IP}
                              \lambda_2\geq6.
\end{equation}
If \(b\leq25\), then \(u,v\leq12\), and
\[
 B(b,r)\leq2r^{12}-1\leq2m^{24}-1<m^{25}
                                           \qquad(m\geq2).
\]
Therefore in the range \(n\geq m^{25}\), every actual
multiplicity obstruction must satisfy
\begin{equation}\label{eq:F014BIP}
                              \lambda_2\geq26.
\end{equation}
Shapes with more than \(d\) rows have zero ambient multiplicity;
one-row shapes have full determinant multiplicity, so neither
exception leaves a gap in these corollaries.
The exclusions hold even when \(d\) is much larger than the
number of rows. They therefore remain additional restrictions
after any independent degree exclusion for \(d(d-1)\leq n\).

\paragraph{Grothendieck dictionary.}
The matrix polynomial ring is the coordinate algebra of the
affine scheme of maps between two locally free sheaves. The
flag minors are exterior-power matrix coefficients; their row
characters and commuting column action are the Cauchy/Schur
functor operations of cote~133, scan~028, and Macdonald, I,
(4.3), p.~63, and Appendix~A, (5.3)--(5.6), pp.~155--156.
The balanced multidegree is a torus-character submodule, with
the grading/comodule correspondence of SGA~3, I, 4.7.3.
Multiplication and division by the displayed common section
give an explicit isomorphism on those characteristic-zero
highest-weight modules. The determinant construction is an
integral identity, whereas irreducible splitting and the
factorial-invertibility argument are used only in characteristic
zero. The induced coherent quotient on the orbit-closure scheme
is what proves equality of the actual multiplicities.
